% ===========================================================================
%  sviluppatore/02-ambiente/04-cartelle.tex
% ===========================================================================
\section{La mappa delle cartelle}\label{sez:cartelle}

Trentuno cartelle nella radice. Raggruppate per mestiere diventano sei
famiglie, e la famiglia dice già quasi tutto su che cosa ci si può aspettare
di trovarci dentro.

\subsection{Il guscio a strati}

Sono le cartelle introdotte con la migrazione all'architettura esagonale
(\vediSez{esagonale}). Sono piccole: la loro dimensione non misura quanto
lavoro fanno, ma quanto poco codice serve a tenere separato il resto.

\begin{center}
\begin{tabular}{@{}L{30mm}R{14mm}L{86mm}@{}}
\toprule
\textbf{Cartella} & \textbf{Righe} & \textbf{Che cosa contiene} \\
\midrule
\rigaTenue
\pkgDominio      & \num{478}  & I tipi puri: \cod{dataclass} congelate, nessuna
                                dipendenza, nemmeno da pandas. \\
\pkgPorte        & \num{239}  & I trentadue contratti. È il file
                                da leggere per primo di tutto il progetto. \\
\rigaTenue
\pkgApplicazione & \num{579}  & I casi d'uso. Conoscono solo le porte. \\
\pkgAdattatori   & \num{2593} & Le implementazioni concrete, e i validatori
                                \simbolo{AST} degli artefatti. \\
\rigaTenue
\pkgComposizione & \num{109}  & Centonove righe, ed è l'unico posto in tutto
                                il sistema dove si sceglie un concreto. \\
\pkgApi          & \num{548}  & L'applicazione FastAPI e i suoi tredici
                                router. \\
\bottomrule
\end{tabular}
\end{center}

\subsection{Il motore}

\begin{center}
\begin{tabular}{@{}L{30mm}R{14mm}L{86mm}@{}}
\toprule
\textbf{Cartella} & \textbf{Righe} & \textbf{Che cosa contiene} \\
\midrule
\rigaTenue
\pkgCore        & \num{7587} & Tipi del dominio quantitativo
                               (\cod{Signal}, \cod{Portfolio}, \cod{Fill}),
                               percorsi, manutenzione. \\
\pkgBacktest    & \num{1567} & Due file. Uno è il motore di simulazione, ed è
                               il punto in cui si decide che cosa è lecito
                               guardare. \\
\pkgEsecuzione  & \num{2428} & I runner: walk-forward, backtest singolo,
                               approcci, gestione dei job. \\
\rigaTenue
\pkgDati        & \num{6251} & Caricamento, universo, categorie, qualità,
                               sorgenti esterne. \\
\pkgReport      & \num{2096} & \simbolo{EnsembleEngine}: la composizione degli
                               stream, i pesi, gli overlay. \\
\rigaTenue
\pkgLive        & \num{4170} & Lo stesso calcolo, collegato a un broker. \\
\bottomrule
\end{tabular}
\end{center}

\subsection{Gli artefatti}

Sono le cinque famiglie di plugin, più tre cartelle che ospitano artefatti
minori. Ogni sottocartella è un artefatto, e il nome della cartella è il nome
canonico --- non quello dichiarato dentro il file.

\begin{center}
\begin{tabular}{@{}L{30mm}R{18mm}L{82mm}@{}}
\toprule
\textbf{Cartella} & \textbf{Installati} & \textbf{Forma del contratto} \\
\midrule
\rigaTenue
\pkgStrategie & \num{9} attive, \num{20} spente
  & \cod{DESCRIPTOR} in \file{strategy.py} \\
\pkgApprocci  & \num{3} attivi, \num{4} dormienti
  & \cod{APPROACH\_DESCRIPTOR} in \file{approach.py} \\
\rigaTenue
\pkgPesi      & \num{14}
  & \cod{META} + \cod{compute} in \file{weight\_mode.py} \\
\pkgOverlay   & \num{7}
  & \cod{META} + \cod{compute} in \file{overlay.py} \\
\rigaTenue
\pkgRegime    & \num{3} regole, \num{3} mercati
  & \cod{META} + \cod{compute}, in due cartelle distinte \\
\pkg{indicators/}, \pkg{metrics/}
  & ---
  & Artefatti minori, stesso schema \\
\bottomrule
\end{tabular}
\end{center}

\begin{pratica}{Accendere e spegnere un artefatto}
Il prefisso \cod{\_} davanti al nome della cartella è l'interruttore:
\file{strategies/\_value/} non viene caricata, \file{strategies/value/} sì.
Rinominare la cartella è l'unico gesto necessario --- non c'è nessun elenco da
aggiornare da nessuna parte. Delle \num{29} cartelle sotto \pkgStrategie,
\num{20} cominciano con il trattino basso e restano fuori.
\end{pratica}

\subsection{L'interfaccia, e il resto}

\pkgDashboard{} è l'applicazione Flask: \num{28494} righe di Python in
\num{36} file, più \num{42072} righe fra template Jinja, JavaScript e CSS. È
l'oggetto più grande del progetto e ha un capitolo suo
(\vediCap{dashboard}).

Restano: \pkg{features/}, \pkg{mms/} e \pkg{volume\_profile/} (calcoli
specialistici che le strategie possono richiedere dichiarandoli nel
descrittore), \pkg{genetic/} (la generazione di strategie per programmazione
genetica), \pkg{validation/} (controlli statistici sui risultati),
\pkg{scripts/} (gli scarichi e le utilità da riga di comando), \pkg{tests/},
\pkg{config/} (il solo \file{quantsys.yaml}) e \pkg{build/} (impacchettamento).

\begin{attenzione}{\pkg{brokers/} è morta}
Nella radice c'è una cartella \pkg{brokers/} di \num{402} righe che duplica
codice vivo in \pkgLive. Nessun modulo la importa. È un residuo: non va
estesa, e ciò che sembra doverci andare va in \file{live/broker.py}.
\end{attenzione}
